\chapter{Analisi Comparativa e Discussione dell'Approccio Ottimo}
\label{chap:capitolo4}

\section{Valutazione quantitativa dei paradigmi analizzati}
\label{sec:valutazione_quantitativa}
% Linee guida:
% - Tabella comparativa dei tempi di calcolo, complessità spaziale e capacità di convergenza verso l'ottimo globale.
% - Risposta degli algoritmi alle metriche di densità del grafo di conflitto specifico dell'Istituto Saffi-Alberti.

Di seguito viene proposta una struttura tabellare per il confronto quantitativo e qualitativo dei diversi approcci analizzati:

\begin{table}[htbp]
\centering
\begin{tabular}{|l|c|c|c|p{4cm}|}
\hline
\textbf{Algoritmo / Paradigma} & \textbf{Tempo di Calcolo} & \textbf{Complessità Spaziale} & \textbf{Ottimalità} & \textbf{Gestione Vincoli Comprensorio (Saffi-Alberti)} \\ \hline
ILP / SAT-Solving              &                           &                               &                     &                                                       \\ \hline
Constraint Programming         &                           &                               &                     &                                                       \\ \hline
Tabu Search                    &                           &                               &                     &                                                       \\ \hline
Simulated Annealing            &                           &                               &                     &                                                       \\ \hline
Algoritmi Genetici             &                           &                               &                     &                                                       \\ \hline
\end{tabular}
\caption{Confronto qualitativo e quantitativo delle metodologie risolutive.}
\label{tab:confronto-metodologie}
\end{table}

\section{Selezione dell'approccio algoritmico e dell'architettura software ideale}
\label{sec:selezione_approccio}
% Linee guida:
% - Motivazioni ingegneristiche a supporto di un approccio ibrido (es. Constraint Programming per la generazione del primo orario ammissibile + Tabu Search per l'ottimizzazione fine dei soft constraints).
% - Analisi di motori di calcolo open-source basati su algoritmi di soddisfacimento di vincoli (es. l'algoritmo euristico ricorsivo di FET) in relazione alla flessibilità di manutenzione del codice.
